% Strict Moodle Heidelberg Beamer version
% Based on the Moodle file beamer-template_en.tex
% Presenter: Junfeng Wei
% Papers:
% 1. Hammond, Vandergheynst, Gribonval, Wavelets on Graphs via Spectral Graph Theory
% 2. Kipf and Welling, Semi-Supervised Classification with Graph Convolutional Networks

\documentclass[aspectratio=1610,xcolor=dvipsnames,t]{beamer}
\usepackage{lmodern}
\usepackage{sansmath}
\usepackage{mathastext}
\usepackage{tikz}
\usetikzlibrary{arrows.meta,positioning,calc}

\usepackage[utf8]{inputenc}
\usepackage{color}
\usepackage{graphicx}
\usepackage{fancybox}

\usepackage{beamerthemesplit}
\usetheme[compress]{Heidelberg}
\definecolor{unirot}{rgb}{0.5976525,0,0}
\usecolortheme[named=unirot]{structure}

\usepackage{amsmath,amssymb,amsfonts}
\usepackage{mathtools}
\usepackage{booktabs}
\usepackage{array}
\usepackage{ragged2e}

\newcolumntype{C}[1]{>{\centering\arraybackslash}p{#1}}
\newcommand{\keyterm}[1]{\textcolor{unirot}{\textbf{#1}}}
\newcommand{\source}[1]{\vspace{0.04cm}{\scriptsize\color{gray} Source: #1}}

\newcommand{\R}{\mathbb{R}}
\newcommand{\C}{\mathbb{C}}
\newcommand{\N}{\mathcal{N}}
\newcommand{\E}{\mathcal{E}}
\newcommand{\V}{\mathcal{V}}
\newcommand{\F}{\mathcal{F}}
\newcommand{\Y}{\mathcal{Y}}
\newcommand{\diag}{\operatorname{diag}}
\newcommand{\softmax}{\operatorname{softmax}}
\newcommand{\ReLU}{\operatorname{ReLU}}
\newcommand{\dist}{\operatorname{dist}}
\newcommand{\supp}{\operatorname{supp}}
\newcommand{\sgwt}{\mathrm{SGWT}}
\newcommand{\gcn}{\mathrm{GCN}}
\newcommand{\what}[1]{\widehat{#1}}
\newcommand{\wt}[1]{\widetilde{#1}}
\newcommand{\inner}[2]{\langle #1,#2\rangle}
\newcommand{\norm}[1]{\left\lVert #1\right\rVert}
\newcommand{\Tbar}{\overline{T}}
\newcommand{\Atil}{\widetilde{A}}
\newcommand{\Dtil}{\widetilde{D}}
\newcommand{\Ahat}{\widehat{A}}
\newcommand{\Ltil}{\widetilde{L}}

\newcommand{\figureplaceholder}[2]{
\begin{center}
\fbox{
\begin{minipage}[c][#2][c]{0.86\textwidth}
\centering
#1
\end{minipage}}
\end{center}
}

\title[From Convolution to GCN]{From Convolution to Graph Convolutional Networks}
\subtitle{Spectral Graph Wavelets, Chebyshev Approximation, and GCNs}
\author[Junfeng Wei, Yizhou Huang]{Junfeng Wei, Yizhou Huang}
\date{\today}
\institute[Uni HD]{
Heidelberg University\\
Institute of Computer Science}

\AtBeginSection[]
{
\begin{frame}<beamer>
\frametitle{Outline}
\tableofcontents[currentsection,hideallsubsections]
\end{frame}
}

\begin{document}

%------------------------------------------------------------
\frame[plain]{\titlepage}

%------------------------------------------------------------
\begin{frame}{Outline}
\begin{enumerate}
  \item Convolution: why translation becomes a problem on graphs
  \item Fourier transform and the convolution theorem
  \item DFT, graph Laplacian, and graph Fourier basis
  \item Spectral Graph Wavelet Transform
  \item Chebyshev wavelet approximation
  \item Graph Convolutional Networks
\end{enumerate}

\medskip
\begin{block}{Core question}
How can convolution-like filtering be defined on irregular graphs without translation?
\end{block}
\end{frame}

\begin{frame}{Logical Roadmap}

\vspace{-0.05cm}

\centering
\resizebox{0.94\textwidth}{!}{%
\begin{tikzpicture}[
    >=Latex,
    mainbox/.style={
        rounded corners=3pt,
        draw=black!45,
        line width=0.7pt,
        fill=gray!6,
        text width=3.10cm,
        minimum height=1.00cm,
        align=center,
        font=\small
    },
    corebox/.style={
        rounded corners=5pt,
        draw=unirot,
        line width=1.15pt,
        fill=unirot!5,
        text width=5.05cm,
        minimum height=1.18cm,
        align=center,
        font=\small
    },
    keybox/.style={
        rounded corners=4pt,
        draw=unirot,
        line width=1.1pt,
        fill=unirot!4,
        text width=3.25cm,
        minimum height=0.95cm,
        align=center,
        font=\small
    },
    smallbox/.style={
        rounded corners=3pt,
        draw=black!45,
        line width=0.7pt,
        fill=gray!6,
        text width=2.95cm,
        minimum height=0.92cm,
        align=center,
        font=\small
    },
    arr/.style={
        ->,
        line width=0.9pt,
        draw=black!65,
        shorten >=2pt,
        shorten <=2pt
    },
    redarr/.style={
        ->,
        line width=0.95pt,
        draw=unirot,
        shorten >=2pt,
        shorten <=2pt
    }
]

% ============================================================
% Top row: classical convolution to graph Fourier basis
% ============================================================

\node[mainbox] (conv) at (0,0)
{Euclidean\\Convolution};

\node[mainbox] (fourier) at (3.90,0)
{Fourier\\Multiplication};

\node[mainbox] (dft) at (7.80,0)
{DFT \& Grid\\Laplacian\\Eigenbasis};

\node[mainbox] (gfb) at (11.70,0)
{Graph Fourier\\Basis};

\draw[arr] (conv.east) -- (fourier.west);
\draw[arr] (fourier.east) -- (dft.west);
\draw[arr] (dft.east) -- (gfb.west);

% ============================================================
% Core node: graph spectral filtering
% Placed above the midpoint of SGWT and Chebyshev
% ============================================================

\node[corebox] (filter) at (4.80,-2.00)
{
{\color{unirot}\textbf{Graph Spectral Filtering}}\\[1mm]
$\displaystyle g(L)=U g(\Lambda)U^\top$
};

% Graph Fourier Basis -> Graph Spectral Filtering
\draw[arr]
(gfb.south) to[out=-90,in=10]
($(filter.east)+(0,0.08)$);

% ============================================================
% Bottom branch: SGWT and GCN route
% ============================================================

\node[keybox] (sgwt) at (2.80,-4.35)
{
{\color{unirot}\textbf{SGWT}}\\[-0.5mm]
$\displaystyle g(tL)$
};

\node[smallbox] (cheb) at (6.80,-4.35)
{Chebyshev\\Approximation};

\node[keybox] (gcn) at (10.55,-4.35)
{
{\color{unirot}\textbf{GCN}}\\[-0.5mm]
Layer
};

% ---------- symmetric clean split from Graph Spectral Filtering ----------
\coordinate (stem) at ($(filter.south)+(0,-0.42)$);

\coordinate (leftcorner)  at ($(stem)+(-2.00,0)$);
\coordinate (rightcorner) at ($(stem)+(2.00,0)$);

\coordinate (sgwttop) at ($(sgwt.north)+(0,0.18)$);
\coordinate (chebtop) at ($(cheb.north)+(0,0.18)$);

% vertical stem
\draw[
    line width=0.95pt,
    draw=unirot
]
(filter.south) -- (stem);

% left branch
\draw[
    line width=0.95pt,
    draw=unirot,
    rounded corners=8pt
]
(stem) -- (leftcorner) -- (sgwttop);

\draw[
    ->,
    line width=0.95pt,
    draw=unirot,
    shorten >=2pt
]
(sgwttop) -- (sgwt.north);

% right branch
\draw[
    line width=0.95pt,
    draw=unirot,
    rounded corners=8pt
]
(stem) -- (rightcorner) -- (chebtop);

\draw[
    ->,
    line width=0.95pt,
    draw=unirot,
    shorten >=2pt
]
(chebtop) -- (cheb.north);

% Chebyshev approximation -> GCN
\draw[arr] (cheb.east) -- (gcn.west);

% ============================================================
% Captions below key nodes
% ============================================================

\node[text=unirot, font=\scriptsize] at (2.80,-5.00)
{$\cdots\cdots\cdots$};

\node[
    align=center,
    font=\scriptsize\bfseries,
    text=unirot
] at (2.80,-5.38)
{multi-scale analysis\\and reconstruction};

\node[text=unirot, font=\scriptsize] at (10.55,-5.00)
{$\cdots\cdots\cdots$};

\node[
    align=center,
    font=\scriptsize\bfseries,
    text=unirot
] at (10.55,-5.38)
{trainable local\\propagation rule};

\end{tikzpicture}
}

\vspace{0.35cm}

\begin{center}
\small
{\Large\textbf{\color{unirot}SGWT}} keeps the multi-scale wavelet viewpoint,\\[1mm]
while {\Large\textbf{\color{unirot}GCN}} simplifies spectral filtering into a trainable propagation layer.
\end{center}

\end{frame}

%============================================================
\section[Convolution]{Convolution}
%============================================================

%------------------------------------------------------------
\begin{frame}{Euclidean Convolution}
For a continuous one-dimensional signal,
\[
(f*h)(a)=\int_{\R} h(a-x)f(x)\,dx.
\]

\begin{itemize}
  \item $f$: input signal.
  \item $h$: filter or kernel.
  \item $a$: output location.
  \item $h(a-x)$: the filter $h$ translated to the output location $a$.
\end{itemize}

\begin{block}{Key point}
Convolution is based on translating one shared filter across a regular domain.
\end{block}
\end{frame}

%------------------------------------------------------------
\begin{frame}{Classical Wavelets}
Classical wavelets are generated from one mother wavelet $\psi$ by translation and scaling:
\[
\psi_{s,a}(x)
=
\frac{1}{s}\psi\left(\frac{x-a}{s}\right).
\tag{SGWT Eq. 1}
\]

For a signal $f$, the wavelet coefficient at scale $s$ and location $a$ is
\[
W_f(s,a)
=
\int_{-\infty}^{\infty}
\frac{1}{s}
\psi^*\left(\frac{x-a}{s}\right)
f(x)\,dx.
\tag{SGWT Eq. 2}
\]

For a fixed scale $s$, we view the transform as an operator:
\[
(T^s f)(a)=W_f(s,a).
\]
\end{frame}

%------------------------------------------------------------
\begin{frame}{Fixed-Scale CWT as Convolution}
Define
\[
\overline{\psi}_s(x)
=
\frac{1}{s}
\psi^*\left(\frac{-x}{s}\right).
\tag{SGWT Eq. 5}
\]

Then, for a fixed scale $s$,
\[
\begin{aligned}
(T^s f)(a)
&=
\int_{-\infty}^{\infty}
\frac{1}{s}
\psi^*\left(\frac{x-a}{s}\right)
f(x)\,dx \\
&=
\int_{-\infty}^{\infty}
\overline{\psi}_s(a-x)f(x)\,dx \\
&=
(\overline{\psi}_s*f)(a).
\end{aligned}
\tag{SGWT Eq. 6}
\]

\begin{block}{Meaning}
At each fixed scale, the classical wavelet transform is a convolution with a scaled wavelet filter.
\end{block}
\end{frame}

%------------------------------------------------------------
\begin{frame}[t]{From Convolution to Spectral Filtering}

\footnotesize
\setlength{\abovedisplayskip}{3pt}
\setlength{\belowdisplayskip}{3pt}
\setlength{\abovedisplayshortskip}{2pt}
\setlength{\belowdisplayshortskip}{2pt}

\begin{alertblock}{Graph difficulty}
Direct translation and scaling are not naturally defined on general graphs.
There is no natural meaning for $a-x$ or $sx$ when $x$ is a graph vertex.
\end{alertblock}

\vspace{-0.4em}

\begin{block}{Key idea}
Move to the Fourier domain, where convolution becomes multiplication.
\end{block}

\vspace{-0.5em}

\[
\widehat{T^s f}(\omega)
=
\widehat{\overline{\psi}_s*f}(\omega)
=
\widehat{\overline{\psi}}_s(\omega)\widehat{f}(\omega).
\tag{\footnotesize SGWT Eq. 7}
\]

\vspace{-0.3em}

Using the scaling property of the Fourier transform,

\vspace{-0.4em}

\[
\widehat{\overline{\psi}}_s(\omega)
=
\widehat{\psi}^{\,*}(s\omega).
\tag{\footnotesize SGWT Eq. 8}
\]

\vspace{-0.4em}

Thus,

\vspace{-0.4em}

\[
(T^s f)(a)
=
\frac{1}{2\pi}
\int_{-\infty}^{\infty}
e^{i\omega a}
\widehat{\psi}^{\,*}(s\omega)
\widehat{f}(\omega)\,d\omega.
\tag{\footnotesize SGWT Eq. 9}
\]

\vspace{-0.2em}

\begin{center}
\keyterm{translation-based convolution}
\quad$\longrightarrow$\quad
\keyterm{spectral filtering}
\end{center}

\end{frame}

%============================================================
\section[Fourier]{Fourier Transform and Convolution Theorem}
%============================================================

\begin{frame}{Fourier Transform}
The Fourier transform and inverse transform are
\[
\widehat f(\omega)
=
\int_{\R}
f(a)e^{-i\omega a}\,da,
\]
\[
f(a)
=
\frac{1}{2\pi}
\int_{\R}
\widehat f(\omega)e^{i\omega a}\,d\omega.
\]

\begin{block}{Notation}
Here $a$ denotes the spatial position, and $\omega$ denotes the Fourier frequency.
\end{block}

\begin{block}{Interpretation}
Fourier transform represents a signal by frequency components rather than by spatial positions.
\end{block}
\end{frame}


\begin{frame}{Convolution Theorem}
For
\[
(f*h)(a)
=
\int_{\R}
f(x)h(a-x)\,dx,
\]
the convolution theorem gives
\[
\widehat{f*h}(\omega)
=
\widehat f(\omega)\widehat h(\omega).
\]

\begin{block}{Consequence}
A convolution operator in the $a$-domain can be represented as pointwise multiplication in the Fourier domain.
\end{block}

\[
\text{convolution in } a\text{-domain}
\quad\Longleftrightarrow\quad
\text{multiplication in } \omega\text{-domain}
\]
\end{frame}


\begin{frame}{Proof Idea of the Convolution Theorem}
\[
\widehat{f*h}(\omega)
=
\int_{\R}
\left(
\int_{\R}
f(x)h(a-x)\,dx
\right)
e^{-i\omega a}\,da.
\]

Let
\[
u=a-x,
\qquad
a=u+x.
\]

Then
\[
e^{-i\omega a}
=
e^{-i\omega(u+x)}
=
e^{-i\omega u}e^{-i\omega x}.
\]

Therefore,
\[
\widehat{f*h}(\omega)
=
\left(
\int_{\R}
f(x)e^{-i\omega x}\,dx
\right)
\left(
\int_{\R}
h(u)e^{-i\omega u}\,du
\right),
\]
\[
\widehat{f*h}(\omega)
=
\widehat f(\omega)\widehat h(\omega).
\]
\end{frame}


\begin{frame}{Fourier Multiplier Form of Wavelets}
From the convolution form of fixed-scale CWT,
\[
(T^s f)(a)
=
(\overline{\psi}_s*f)(a).
\]

Taking the Fourier transform gives
\[
\widehat{T^s f}(\omega)
=
\widehat{\overline{\psi}_s}(\omega)\widehat f(\omega).
\tag{SGWT Eq. 7}
\]

Using the scaling property,
\[
\widehat{\overline{\psi}_s}(\omega)
=
\widehat\psi^{\,*}(s\omega).
\tag{SGWT Eq. 8}
\]

Therefore,
\[
(T^s f)(a)
=
\frac{1}{2\pi}
\int_{-\infty}^{\infty}
e^{i\omega a}
\widehat\psi^{\,*}(s\omega)
\widehat f(\omega)\,d\omega.
\tag{SGWT Eq. 9}
\]
\end{frame}


\begin{frame}{Why This Matters for Graphs}
\begin{block}{Classical wavelets}
Scaling can be expressed in the Fourier domain as
\[
\widehat\psi^{\,*}(s\omega).
\]
\end{block}

\begin{block}{Graph wavelets}
Replace Fourier frequencies $\omega$ by graph Laplacian eigenvalues $\lambda_\ell$:
\[
\omega
\quad\longrightarrow\quad
\lambda_\ell.
\]
\end{block}

\[
\widehat\psi^{\,*}(s\omega)
\quad\leadsto\quad
g(t\lambda_\ell).
\]

\begin{center}
\keyterm{This is the conceptual bridge from CWT to SGWT.}
\end{center}
\end{frame}

%============================================================
\section[DFT and Graph Fourier]{DFT, Graph Laplacian, and Graph Fourier Basis}
%============================================================

%------------------------------------------------------------
\begin{frame}[t]{DFT on Finite Regular Grids}

\setlength{\abovedisplayskip}{3pt}
\setlength{\belowdisplayskip}{3pt}
\setlength{\abovedisplayshortskip}{2pt}
\setlength{\belowdisplayshortskip}{2pt}

For a finite discrete signal $x[n]$, the DFT is
\[
X[k]
=
\sum_{n=0}^{N-1}
x[n]e^{-2\pi i kn/N}.
\]

\vspace{-0.5em}

The inverse DFT is
\[
x[n]
=
\frac{1}{N}
\sum_{k=0}^{N-1}
X[k]e^{2\pi i kn/N}.
\]

\vspace{-0.5em}

The DFT basis functions are
\[
\chi_k[n]
=
e^{2\pi i kn/N}.
\]

\vspace{-0.4em}

\begin{block}{Key point}
DFT gives a Fourier basis for finite regular grids:
\[
f(a)\leftrightarrow x[n],
\quad
\widehat f(\omega)\leftrightarrow X[k],
\quad
e^{i\omega a}\leftrightarrow e^{2\pi i kn/N}.
\]
This will serve as the reference case when we replace regular-grid Fourier modes by graph Laplacian eigenvectors.
\end{block}

\end{frame}

%------------------------------------------------------------
\begin{frame}{DFT Convolution Theorem}

For convolution,
\[
(x*h)[n]
=
\sum_{m=0}^{N-1}
x[m]h[(n-m)\bmod N],
\]
\[
(f*h)(a)
\leftrightarrow
(x*h)[n].
\]

we have
\[
\mathrm{DFT}(x*h)[k]
=
X[k]H[k],
\]
\[
\widehat{f*h}(\omega)
=
\widehat f(\omega)\widehat h(\omega)
\quad\leftrightarrow\quad
\mathrm{DFT}(x*h)[k]
=
X[k]H[k].
\]

\begin{block}{Key point}
The DFT keeps the same convolution principle as the continuous Fourier transform: convolution in the original domain becomes multiplication in the frequency domain.
\end{block}

\end{frame}

%------------------------------------------------------------
\begin{frame}{Graph Laplacian and Spectrum}
For a weighted undirected graph,
\[
L=D-A.
\]

Since $L$ is real symmetric, it has an orthonormal eigenbasis:
\[
L\chi_\ell
=
\lambda_\ell\chi_\ell,
\]
\[
0=\lambda_0
\leq
\lambda_1
\leq
\cdots
\leq
\lambda_{N-1}.
\]

For the unnormalized Laplacian,
\[
f^TLf
=
\frac{1}{2}
\sum_{i,j}
A_{ij}(f_i-f_j)^2.
\]

\begin{center}
\keyterm{Small eigenvalues correspond to smooth graph signals.}
\end{center}
\end{frame}

%------------------------------------------------------------
\begin{frame}{Graph Fourier Transform}
The graph Fourier transform is
\[
\widehat f(\ell)
=
\langle \chi_\ell,f\rangle
=
\sum_{n=0}^{N-1}
\chi_\ell^*(n)f(n).
\tag{SGWT Eq. 18}
\]

The inverse graph Fourier transform is
\[
f(n)
=
\sum_{\ell=0}^{N-1}
\widehat f(\ell)\chi_\ell(n).
\tag{SGWT Eq. 19}
\]

The Parseval relation is
\[
\langle f,h\rangle
=
\langle \widehat f,\widehat h\rangle.
\tag{SGWT Eq. 20}
\]

\keyterm{The inner product of two graph signals is the same in the vertex domain and in the graph Fourier domain.}
\end{frame}

%------------------------------------------------------------
\begin{frame}{From Graph Fourier to Graph Filtering}
\begin{block}{Fourier multiplier}
\[
\widehat{Tf}(\omega)
=
m(\omega)\widehat f(\omega).
\]
\end{block}

\begin{block}{Graph spectral multiplier}
\[
\widehat{T_g f}(\ell)
=
g(\lambda_\ell)\widehat f(\ell).
\]
\end{block}

\begin{center}
\keyterm{This gives a way to define graph filters without graph translation.}
\end{center}
\end{frame}

%============================================================
\section[SGWT]{Spectral Graph Wavelet Transform}
%============================================================
\begin{frame}{Spectral Graph Filter}

For a spectral filter function $g$, define the graph spectral operator by
\[
\widehat{T_g f}(\ell)
=
g(\lambda_\ell)\widehat f(\ell).
\tag{SGWT Eq.~21}
\]

In the vertex domain, this becomes
\[
(T_g f)(m)
=
\sum_{\ell=0}^{N-1}
g(\lambda_\ell)\widehat f(\ell)\chi_\ell(m).
\tag{SGWT Eq.~22}
\]

Equivalently, if $L=U\Lambda U^*$, then
\[
T_g
=
g(L)
=
U g(\Lambda)U^*.
\]

\vspace{0.4em}

\end{frame}

\begin{frame}{Scaled Spectral Filter}

The scaled graph wavelet operator is
\[
T_g^t = g(tL).
\]

With eigendecomposition $L = U\Lambda U^*$,
\[
g(tL)
=
U g(t\Lambda) U^*.
\]

For each graph frequency $\lambda_\ell$,
\[
g(tL)\chi_\ell
=
g(t\lambda_\ell)\chi_\ell.
\]

\begin{block}{Scale in SGWT}
The scale parameter $t$ rescales graph frequencies:
\[
g(\lambda_\ell)
\quad\longrightarrow\quad
g(t\lambda_\ell).
\]
It does not rely on vertex-domain translation.
\end{block}

\end{frame}

\begin{frame}{Graph Wavelets}

The graph wavelet centered at vertex $n$ and scale $t$ is
\[
\psi_{t,n}
=
T_g^t\delta_n
=
g(tL)\delta_n.
\tag{SGWT Eq.~23}
\]

The value at vertex $m$ is
\[
\psi_{t,n}(m)
=
\sum_{\ell=0}^{N-1}
g(t\lambda_\ell)\chi_\ell^*(n)\chi_\ell(m).
\tag{SGWT Eq.~24}
\]

\begin{block}{Graph wavelet}
A graph wavelet is obtained by applying the scaled spectral operator $g(tL)$
to an impulse $\delta_n$ at vertex $n$.
\end{block}

\end{frame}

\begin{frame}{Wavelet Coefficients}

The SGWT coefficient is
\[
W_f(t,n)
=
\inner{\psi_{t,n}}{f}.
\tag{SGWT Eq.~25}
\]

Equivalently,
\[
W_f(t,n)
=
(T_g^t f)(n)
=
\sum_{\ell=0}^{N-1}
g(t\lambda_\ell)\widehat f(\ell)\chi_\ell(n).
\tag{SGWT Eq.~26}
\]

\begin{block}{Operational form}
The coefficient $W_f(t,n)$ is the value at vertex $n$ after filtering $f$ by $g(tL)$.
\end{block}

\end{frame}

\begin{frame}{Two Key Properties of SGWT}

After defining the SGWT coefficients, we need to check two key properties.

\begin{block}{1. Locality}
Are graph wavelets concentrated around their center vertex?
\end{block}

\begin{block}{2. Stability}
Do the coefficients give a stable representation of the signal?
\end{block}

\end{frame}

\begin{frame}{Localization by Graph Distance}

Let $d_G(m,n)$ be the graph distance between vertices $m$ and $n$.

\medskip

A key consequence of the path-counting lemmas is that powers of the graph Laplacian are localized in graph distance:
\[
d_G(m,n)>s
\quad\Rightarrow\quad
(L^s)_{m,n}=0.
\]

\begin{itemize}
  \item One multiplication by $L$ propagates information by at most one hop.
  \item $L^s$ propagates information by at most $s$ hops.
  \item Polynomial filters in $L$ are localized in graph distance.
\end{itemize}

\begin{block}{Key idea}
Matrix powers encode graph neighborhoods: $L^s$ cannot connect vertices farther than $s$ hops.
\end{block}

\end{frame}

\begin{frame}{Taylor Structure of the Kernel}

If $g$ has a zero of order $K$ at the origin, then for small $t$,
\[
g(t\lambda)
\approx
\frac{g^{(K)}(0)}{K!}(t\lambda)^K.
\]

Therefore, at the operator level,
\[
g(tL)
\approx
\frac{g^{(K)}(0)}{K!}t^K L^K.
\]

\begin{block}{Connection to localization}
Since $L^K\delta_n$ cannot reach vertices farther than $K$ hops,
the leading term of the graph wavelet is localized around vertex $n$.
\end{block}

\end{frame}

\begin{frame}{Theorem 5.7: Small-Scale Localization}

Let $g$ satisfy the Taylor assumption above.
If
\[
d_G(m,n)>K,
\]
then for sufficiently small $t$,
\[
\frac{|\psi_{t,n}(m)|}{\|\psi_{t,n}\|_2}
\leq
D_{m,n}t.
\tag{SGWT Eq.~42}
\]

\begin{block}{Meaning}
At small scales, graph wavelets are localized around their center vertex.
Far-away vertices have only a small relative response.
\end{block}

\end{frame}

\begin{frame}{Proof Idea: Theorem 5.7}

Use the Taylor expansion of the kernel:
\[
g(tL)
=
\frac{g^{(K)}(0)}{K!}t^K L^K
+
O(t^{K+1}).
\]

Apply it to the impulse $\delta_n$:
\[
\psi_{t,n}
=
g(tL)\delta_n
=
\frac{g^{(K)}(0)}{K!}t^K L^K\delta_n
+
O(t^{K+1}).
\]

If $d_G(m,n)>K$, then
\[
(L^K)_{m,n}=0.
\]

Hence the leading term vanishes at vertex $m$:
\[
|\psi_{t,n}(m)|=O(t^{K+1}),
\qquad
\|\psi_{t,n}\|_2=O(t^K).
\]

Therefore,
\[
\frac{|\psi_{t,n}(m)|}{\|\psi_{t,n}\|_2}
=
O(t).
\]

\begin{center}
\keyterm{At small scales, the wavelet is concentrated near its center vertex.}
\end{center}

\end{frame}

\begin{frame}{Lemma 5.1: Continuous SGWT Inverse}

Assume
\[
C_g=\int_0^\infty \frac{g(x)^2}{x}\,dx<\infty.
\tag{SGWT Eq.~27}
\]

Then the continuous family of graph wavelets reconstructs
\[
f^\#=f-\inner{\chi_0}{f}\chi_0.
\]

The reconstruction formula is
\[
\frac{1}{C_g}\sum_{n=1}^{N}\int_0^\infty
W_f(t,n)\psi_{t,n}(m)\frac{dt}{t}
=f^\#(m).
\tag{SGWT Eq.~28}
\]

\end{frame}

\begin{frame}{Proof Idea: Lemma 5.1}

Expand both $W_f(t,n)$ and $\psi_{t,n}(m)$ in the graph Fourier basis:
\[
W_f(t,n)
=
\sum_{\ell=0}^{N-1}
g(t\lambda_\ell)\widehat f(\ell)\chi_\ell(n),
\]
\[
\psi_{t,n}(m)
=
\sum_{\ell'=0}^{N-1}
g(t\lambda_{\ell'})\chi_{\ell'}^*(n)\chi_{\ell'}(m).
\]

Sum over the center vertex $n$:
\[
\sum_{n=1}^{N}\chi_\ell(n)\chi_{\ell'}^*(n)=\delta_{\ell,\ell'}.
\]

Then integrate over scales:
\[
\int_0^\infty g(t\lambda_\ell)^2\frac{dt}{t}=C_g
\quad(\lambda_\ell>0).
\]

\begin{center}
\keyterm{Only the zero-frequency term is missing, because $g(0)=0$.}
\end{center}

\end{frame}

\begin{frame}{Scaling Functions}

Wavelet filters usually satisfy $g(0)=0$, so they do not capture the zero-frequency component.

\medskip

To cover the low-frequency part, SGWT introduces a low-pass kernel $h$.
The scaling function centered at vertex $n$ is
\[
\phi_n
=
h(L)\delta_n.
\]

The scaling coefficient is
\[
S_f(n)
=
\inner{\phi_n}{f}.
\]

\begin{block}{Low-frequency coverage}
The wavelet kernels $g(t\lambda_\ell)$ capture band-pass components,
while the scaling kernel $h(\lambda_\ell)$ covers the low-frequency part near $\lambda=0$.
\end{block}

\end{frame}

\begin{frame}{Spectral Coverage Criterion}

For stable reconstruction, the scaling filter and all wavelet filter should cover the graph spectrum.

For each graph frequency $\lambda_\ell$, define
\[
G(\lambda_\ell)
=
|h(\lambda_\ell)|^2
+
\sum_{j=1}^{J}
|g(t_j\lambda_\ell)|^2.
\]

A stable SGWT requires
\[
0 < A
\leq
G(\lambda_\ell)
\leq
B < \infty,
\qquad
\ell=0,\dots,N-1.
\]

\begin{block}{Meaning}
No graph frequency should be missed.  
The scaling kernel $h$ covers the low-frequency part, especially near $\lambda=0$,
while the scaled wavelet kernels $g(t_j\lambda_\ell)$ cover the remaining frequencies.
\end{block}

\end{frame}

\begin{frame}{Stability of SGWT Coefficients}

The spectral coverage condition implies the frame inequality
\[
A\|f\|^2
\leq
\sum_n |S_f(n)|^2
+
\sum_{j=1}^{J}\sum_n |W_f(t_j,n)|^2
\leq
B\|f\|^2.
\]

Here, the middle term is the total energy of all SGWT coefficients.

\begin{block}{Interpretation}
The lower bound means that no signal energy is lost.
The upper bound means that the coefficients do not over-amplify the signal.
\end{block}

\end{frame}

%============================================================
\section[Approximation]{Chebyshev Wavelet Approximation}
%============================================================

\begin{frame}{Why Approximation Is Needed}
PROBLEM: Direct spectral filtering is computationally expensive.

Specifically, eigen-decomposition of the wavelet generating function
$$g(tL)f=Ug(t\Lambda)U^*f$$
scales poorly at $O(N^3)$. For large graphs, computing $U$ is expensive.

\medskip
SOLUTION: Instead of finding the exact, filter, approximate using some polynomial functions instead: 
$$g(tL)f\approx p(L)f$$
where
$$p(L)=\sum_{i=0}^M c_i(L)^i$$

\medskip
QUESTION: Will this shortcut break the function?
\end{frame}

\begin{frame}{Lemma 6.1: Approximation Error}
Let approximate $p(x)$ of $g(tx)$ over $[0,\lambda_{\max}]$. Define spectral error as
$$B:=\sup_{x\in[0,\lambda_{\max}]}|g(tx)-p(x)|.$$

Define the approximate coefficient at node $n$ as
$$\widetilde{W}_f(t,n):=(p(L)f)_n.$$

Then
\[
|W_f(t,n)-\widetilde W_f(t,n)|\leq B\|f\|.
\tag{SGWT Eq. 53}
\]

\begin{center}
\keyterm{The error at any node is bounded by $B$, $||f||$ is a scalar.}

\medskip
Blackboard presentation for proof.
\end{center}
\end{frame}

\begin{frame}{Chebyshev Polynomials}
One potential candidate for $p(L)$ are the Chebyshev polynomials:
\[
T_0(y)=1,
\qquad
T_1(y)=y,
\]
\[
T_k(y)=2yT_{k-1}(y)-T_{k-2}(y).
\]

They also satisfy
\[
T_k(y)=\cos(k\arccos y),
\qquad y\in[-1,1].
\]

\begin{block}{Reason for use}
Chebyshev polynomials are numerically stable on $[-1,1]$ and can be evaluated recursively.
\end{block}
\end{frame}

\begin{frame}{Shifted Chebyshev Approximation}
Map $[0,\lambda_{\max}]$ to $[-1,1]$ by
\[
a=\frac{\lambda_{\max}}{2},
\qquad
y=\frac{x-a}{a}.
\]

Recall that
$$T_k(y)=2yT_{k-1}(y)-T_{k-2}(y),$$
the shifted matrix-vector equivalent is
$$\overline{T}_k(L)f=\frac{2(L-aI)}{a}\overline{T}_{k-1}(L)f-\overline{T}_{k-2}(L)f.$$

Therefore,
\[
g(t_jL)f
\approx
\frac{1}{2}c_{j,0}f+
\sum_{k=1}^{M_j}c_{j,k}\overline T_k(L)f.
\]
\end{frame}

\begin{frame}{Fast Recurrence}
The shifted recurrence is
\begin{align*}
    \overline T_k(L)f&=\frac{2(L-aI)}{a}\overline T_{k-1}(L)f-\overline T_{k-2}(L)f\\
    &=\frac{2}{a}(L-aI)v_{k-1}-v_{k-2}
\end{align*}

\begin{itemize}
  \item $M$ matrix-vector multiplication and subtractions.
  \item For $p(L):=\sum_{k=0}^Kc_k(L)^k$, filtering costs approximately $O(K|\E|)$.
\end{itemize}
\end{frame}

\begin{frame}{Reconstruction with the approximated SGWT}
SGWT creates a $N\times(J+1)$ dimensional "overcomplete frame":
\[
Wf=
\left(
T_hf,\,
T_g^{t_1}f,\,
\ldots,\,
T_g^{t_J}f
\right).
\]

A pseudoinverse is needed:
\[
W^\dagger=(W^*W)^{-1}W^*.
\]

This requires the adjoint of $W$:
\[
W^*\eta
=
T_h\eta_0+\sum_{j=1}^{J}T_g^{t_j}\eta_j.
\]

\begin{itemize}
    \item $W^*$ applies scaling filter to low-pass coefficients and wavlet filters to band-pass coefficients.
    \item weighted sum to project multi-sclae wavelets back to the vertices.
\end{itemize}
\end{frame}

\begin{frame}{Reconstruction with the approximated SGWT (Continued)}

Pseudoinverse reconstruction:
$$f=(W^*W)^{-1}W^*\eta.$$

Let $b=W^*\eta$,
find $f$ s.t.
$$(W^*W)f=b.$$

\medskip

Calculate using "Conjugate Gradients" algorithm:

\begin{itemize}
    \item Compute $f_{temp}=Wf$ using Chebyshev recurrence.
    \item Compute $f_{final}=W^*f_{temp}$ using Chebyshev recurrence.
\end{itemize}
\end{frame}

%============================================================
\section[GCN]{Graph Convolutional Networks}
%============================================================

\begin{frame}{Semi-Supervised Node Classification}
Inputs:
\[
A\in\R^{N\times N},
\qquad
X\in\R^{N\times C},
\qquad
\Y_L.
\]

Model output:
\[
Z=f(X,A)\in\R^{N\times F}.
\]

The loss is evaluated only on labeled nodes,
\[
\Y_L\subset\{1,\ldots,N\}.
\]

\begin{block}{GCN model class}
Standard NN only takes $X$ as input and output $Z=f(X)$. GCN introduces graph structure $A$ directly into forward pass.
\end{block}
\end{frame}

\begin{frame}{Laplacian Regularization}
Traditionally, the loss function is
\[
\mathcal{L}=\mathcal{L}_0+\lambda\mathcal{L}_{\mathrm{reg}}.
\tag{GCN Eq. 1}
\]
where
\[
\mathcal{L}_{\mathrm{reg}}
=
\sum_{i,j}A_{ij}\|f(X_i)-f(X_j)\|^2
=f(X)^T\Delta f(X).
\]

\begin{alertblock}{Limitation}
Connected nodes must have the same label, forcing output to bet smoothed across the graph.
\end{alertblock}
\end{frame}

\begin{frame}{GCN Spectral Convolution}
Normalized graph Laplacian is needed:
\[
L=I_N-D^{-1/2}AD^{-1/2}=U\Lambda U^T.
\]
Eigenvalues of $L$ bounded between $[0,2]$.

GCN defines spectral convolution by
\[
g_\theta*x=Ug_\theta U^Tx.
\tag{GCN Eq. 3}
\]

\begin{itemize}
  \item $U^Tx$: graph Fourier transform.
  \item $g_\theta$: learnable spectral filter.
  \item $U$: inverse graph Fourier transform.
\end{itemize}
\end{frame}

\begin{frame}{Chebyshev Approximation in GCN}
Scale eigenvalues to Chebyshev stable range
\[
\widetilde{\Lambda}=\frac{2}{\lambda_{\max}}\Lambda-I_N.
\]

Chebyshev polynomial approximation:
\[
g_{\theta'}(\Lambda)
\approx
\sum_{k=0}^{K}\theta_k'T_k(\widetilde{\Lambda}),
\tag{GCN Eq. 4}
\]

Thus,
\[
g_{\theta'}*x
\approx
\sum_{k=0}^{K}\theta_k'T_k(\widetilde L)x,
\tag{GCN Eq. 5}
\]
with
\[
\widetilde L=\frac{2}{\lambda_{\max}}L-I_N.
\]
\end{frame}

\begin{frame}{First-Order Approximation}
Set $K=1$ and approximate $\lambda_{\max}\approx 2$:
\[
g_{\theta'}*x
\approx
\theta_0'x+\theta_1'(L-I_N)x.
\]

Since
\[
L=I_N-D^{-1/2}AD^{-1/2},
\]
\[
g_{\theta'}*x
\approx
\theta_0'x-\theta_1'D^{-1/2}AD^{-1/2}x.
\tag{GCN Eq. 6}
\]

To prevent overfitting, fix
\[
\theta:=\theta_0'=-\theta_1'.
\]

Then
\[
g_\theta*x
\approx
\theta\left(I_N+D^{-1/2}AD^{-1/2}\right)x.
\tag{GCN Eq. 7}
\]

\begin{alertblock}{Problem}
The operator $I_N+D^{-1/2}AD^{-1/2}$ has eigenvalues in $[0,2]$, which can cause numerical instability when layers are stacked.
\end{alertblock}
\end{frame}

\begin{frame}{Renormalization Trick}
Add self-loops:
\[
\widetilde A=A+I_N,
\qquad
\widetilde D_{ii}=\sum_j \widetilde A_{ij}.
\]

The renormalized propagation model is
\[
Z=\widetilde D^{-1/2}\widetilde A\widetilde D^{-1/2}X\Theta.
\tag{GCN Eq. 8}
\]

Define
\[
\widehat A=\widetilde D^{-1/2}\widetilde A\widetilde D^{-1/2}.
\]

\begin{center}
\keyterm{Operator's eigenvalues are now bounded within $[-1,1]$.}
\end{center}
\end{frame}

\begin{frame}{Layer-Wise GCN Propagation (Macro)}
Stacking the renormalized propagation with trainable weights and nonlinearities gives
\[
H^{(l+1)}
=
\sigma\left(\widehat A H^{(l)}W^{(l)}\right).
\tag{GCN Eq. 2}
\]

\begin{itemize}
  \item $H^{(0)}=X$.
  \item $\widehat A H^{(l)}$: normalized neighborhood aggregation.
  \item $W^{(l)}$: learnable feature transformation.
  \item $\sigma$: nonlinearity such as ReLU.
\end{itemize}
\end{frame}

\begin{frame}{Neighborhood Aggregation View (Micro)}
GCN can be written node-wise as
\[
h_i^{(l+1)}
=
\sigma\left(
\sum_{j\in\mathcal{N}_i}
\frac{1}{c_{ij}}h_j^{(l)}W^{(l)}
\right).
\tag{GCN Eq. 12}
\]

For symmetric normalization,
\[
c_{ij}=\sqrt{d_i d_j}.
\]

\begin{block}{Connection}
This is the vertex-domain aggregation form of the first-order spectral approximation.
\end{block}
\end{frame}

\begin{frame}{Two-Layer GCN}
For semi-supervised classification, the two-layer model is
\[
Z=\softmax\left(
\widehat A\,
\sigma\left(\widehat A XW^{(0)}\right)
W^{(1)}
\right).
\tag{GCN Eq. 9}
\]

\begin{itemize}
  \item First propagation: feature smoothing and hidden representation.
  \item Second propagation: class-level output representation.
  \item Two layers aggregate information from up to two-hop neighborhoods.
\end{itemize}

The cross-entropy loss over labeled nodes is
\[
\mathcal{L}
=
-
\sum_{l\in\mathcal{Y}_L}
\sum_{f=1}^{F}
Y_{lf}\ln Z_{lf}.
\tag{GCN Eq. 10}
\]

\begin{block}{Semi-supervised setting}
Unlabeled nodes are not included in the loss, but they participate in message propagation through $\widehat A$.
\end{block}
\end{frame}

\begin{frame}{GCN Main Results}
\begin{table}
\scriptsize
\centering
\resizebox{\textwidth}{!}{
\begin{tabular}{lrrrr}
\toprule
Method & Citeseer & Cora & Pubmed & NELL \\
\midrule
ManiReg & 60.1 & 59.5 & 70.7 & 21.8 \\
SemiEmb & 59.6 & 59.0 & 71.1 & 26.7 \\
LP & 45.3 & 68.0 & 63.0 & 26.5 \\
DeepWalk & 43.2 & 67.2 & 65.3 & 58.1 \\
ICA & 69.1 & 75.1 & 73.9 & 23.1 \\
Planetoid* & 64.7 (26s) & 75.7 (13s) & 77.2 (25s) & 61.9 (185s) \\
GCN & 70.3 (7s) & 81.5 (4s) & 79.0 (38s) & 66.0 (48s) \\
GCN random splits & $67.9\pm0.5$ & $80.1\pm0.5$ & $78.9\pm0.7$ & $58.4\pm1.7$ \\
\bottomrule
\end{tabular}}
\end{table}

\begin{center}
\keyterm{GCN improves accuracy and is faster than previous best methods like Planetoid.}
\end{center}
\source{GCN Table 2.}
\end{frame}

\begin{frame}{GCN Propagation Rule Ablation}
\begin{table}
\scriptsize
\centering
\resizebox{\textwidth}{!}{
\begin{tabular}{lrrr}
\toprule
Propagation model & Citeseer & Cora & Pubmed \\
\midrule
Chebyshev filter, $K=3$ & 69.8 & 79.5 & 74.4 \\
Chebyshev filter, $K=2$ & 69.6 & 81.2 & 73.8 \\
First-order model, Eq. 6 & 68.3 & 80.0 & 77.5 \\
Single parameter, Eq. 7 & 69.3 & 79.2 & 77.4 \\
Renormalization trick, Eq. 8 & 70.3 & 81.5 & 79.0 \\
First-order term only & 68.7 & 80.5 & 77.8 \\
Multi-layer perceptron & 46.5 & 55.1 & 71.4 \\
\bottomrule
\end{tabular}}
\end{table}

\begin{center}
\keyterm{Ablation shows renormalisation is the best performing variant.}
\end{center}
\source{GCN Table 3.}
\end{frame}

\begin{frame}{GCN Runtime and Depth Evidence}
\begin{figure}
        \centering
        \includegraphics[width=0.6\textwidth]{wallclock time.png}
\end{figure}

\begin{block}{Runtime scaling, Figure 2}
On random graphs, Chebyshev approximation reduces computation to linear time.
\end{block}
\end{frame}

\begin{frame}{GCN Runtime and Depth Evidence (Continued)}
\begin{figure}
        \centering
        \includegraphics[width=0.9\textwidth]{model depth.png}
\end{figure}

\begin{center}
\keyterm{Peak of curves are highest at 2 to 3 layers. Deeper layers struggle to train.}
\end{center}
\end{frame}

\begin{frame}{Limitations of GCN}
\begin{itemize}
  \item Full-batch training has memory cost growing linearly with graph size.
  \item Mini-batch training is nontrivial because a $K$-layer GCN depends on $K$-hop neighborhoods.
  \item The original model is limited to undirected graphs without native edge features.
  \item The approximation assumes locality and equal importance of self-connections and neighbor edges.
\end{itemize}

\begin{block}{Suggested extension}
\[
\widetilde A=A+\lambda I_N
\tag{GCN Eq. 11}
\]
can introduce a learnable or tunable trade-off between self-information and neighbor information.
\end{block}
\end{frame}

%============================================================
\section[End]{Summary and Conclusion}
%============================================================

\begin{frame}{SGWT vs. GCN}
\begin{columns}[T]
\begin{column}{0.48\textwidth}
\textbf{SGWT}
\begin{itemize}
  \item Signal analysis on graphs.
  \item Multi-scale wavelets.
  \item Reconstruction and frame theory.
  \item Core operator:
  \[
  g(tL)
  \]
\end{itemize}
\end{column}

\begin{column}{0.48\textwidth}
\textbf{GCN}
\begin{itemize}
  \item Node classification.
  \item Learnable feature propagation.
  \item First-order local approximation.
  \item Core layer:
  \[
  H^{(l+1)}=\sigma(\widehat A H^{(l)}W^{(l)})
  \]
\end{itemize}
\end{column}
\end{columns}
\end{frame}

\begin{frame}{References}
\small
\begin{itemize}
  \item D. K. Hammond, P. Vandergheynst, and R. Gribonval.\newline
  \textit{Wavelets on Graphs via Spectral Graph Theory}.\newline
  Applied and Computational Harmonic Analysis, 2011.

  \item T. N. Kipf and M. Welling.\newline
  \textit{Semi-Supervised Classification with Graph Convolutional Networks}.\newline
  International Conference on Learning Representations, 2017.

  \item M. Defferrard, X. Bresson, and P. Vandergheynst.\newline
  \textit{Convolutional Neural Networks on Graphs with Fast Localized Spectral Filtering}.\newline
  NeurIPS, 2016.
\end{itemize}
\end{frame}

\begin{frame}{Questions}
\vfill
\begin{center}
{\Huge\color{unirot}\textbf{Thank you.}}\\[0.4cm]
{\Large Questions and Discussion}
\end{center}
\vfill
\end{frame}

\end{document}



